%% Методические указания
%% Лаборатория математики ФН1
\documentclass[12pt, a4paper]{article}
\usepackage[T2A]{fontenc}
\usepackage[utf8]{inputenc}
\usepackage[russian]{babel}
\usepackage[left=25mm, right=20mm, top=20mm, bottom=20mm]{geometry}
\usepackage{amsmath, amssymb}
\usepackage{enumitem}
\pagestyle{plain}
\setlength{\parindent}{1.25cm}
\begin{document}
\begin{center}
  {\large\bfseries Методические указания}

  \vspace{0.3em}
  {\large Случайные величины: как не ошибиться в модели}

  \vspace{0.5em}
  {\small Теория вероятностей · 2 курс\\ Лаборатория математики ФН1}
\end{center}
\vspace{1em}

\section*{С чего начинать задачу}

Большинство ошибок в теорвере~--- не арифметические, а модельные.
Прежде чем считать, ответьте на три вопроса.

\begin{enumerate}[nosep]
  \item \textbf{Что является элементарным исходом?} Выпишите пространство $\Omega$ явно.
  \item \textbf{Равновероятны ли исходы?} Классическое определение вероятности
        применимо только к равновозможным исходам.
  \item \textbf{Независимы ли события?} Независимость нужно обосновывать условием задачи,
        а не предполагать по умолчанию.
\end{enumerate}

\section*{Дискретная величина}

\begin{itemize}[nosep]
  \item Ряд распределения: все значения и их вероятности; сумма вероятностей равна единице~---
        это обязательная проверка.
  \item Функция распределения $F(x) = P(X < x)$ ступенчатая, непрерывна слева,
        и это надо отразить на графике выколотыми точками.
  \item $M[X] = \sum x_i p_i$, \quad $D[X] = M[X^2] - (M[X])^2$.
\end{itemize}

\section*{Непрерывная величина}

\begin{itemize}[nosep]
  \item Плотность неотрицательна, и её интеграл по всей прямой равен единице~---
        отсюда находится неизвестный коэффициент.
  \item $F(x) = \int_{-\infty}^{x} f(t)\,dt$; проверьте, что $F$ непрерывна и не убывает.
  \item $M[X] = \int x f(x)\,dx$, \quad $D[X] = \int x^2 f(x)\,dx - (M[X])^2$.
\end{itemize}

\section*{Типичные ошибки}

\begin{itemize}[nosep]
  \item Применяют формулу Бернулли к зависимым испытаниям.
  \item Забывают проверить нормировку и получают вероятности больше единицы.
  \item Путают плотность и функцию распределения при вычислении $M[X]$.
  \item В формуле полной вероятности берут гипотезы, не образующие полной группы.
\end{itemize}

\end{document}
